\subsection{代数张量积上的模与多模}

定义 2. 设 E 是一个（含单位元的）A-代数。左（相应地，右）E-模就是 E 的底环上的一个左（相应地，右）模。

除非另有说明，本号中所考虑的模和多重模均为左模与左多重模。

若 M 是一个 E-模，则同态 \(\eta: A \to E\) （ \(\S 1\) , 第 4 号）在 M 上定义了一个 A-模结构，称为 M 上 E-模结构之下的结构；对 \(\alpha \in A\) , \(s \in E\) , \(x \in M\) ,

\[
\alpha (s x) = s (\alpha x) = (\alpha s) x,\tag{7}
\]

从而对所有 \(s \in E\) ， M 的位似 \(h_{s}: x \mapsto sx\) 是底层 A-模结构的一个自同态。反之，在 M 上给定一个 E-模结构，等价于在 M 上给定一个 A-模结构以及一个 A-代数同态 \(s \mapsto h_{s}\) （从 E 到 \(\operatorname{End}_{\mathbf{A}}(\mathbf{M})\) ）.

定义 3. 设 E 和 F 是两个（含单位元的）A-代数，且 M 是一个同时具有 E-模结构和 F-模结构的集合。若 M 满足以下条件，则称 M 为代数 E 和 F 上的（左）双模：

\begin{enumerate}
\def\labelenumi{(\arabic{enumi})}
\item
  M 是 E 和 F 的底环上的双模（II, §1, 第 14 号）；
\item
  M 上承载 E-模与 F-模结构的两个 A-模结构是相同的。
\end{enumerate}

后一条件是说，若 \(e\) 与 \(e'\) 分别是 \(\mathbf{E}\) 与 \(\mathbf{F}\) 的单位元，则

\[
(\alpha e) x = (\alpha e ^ {\prime}) x \quad \text { for } \alpha \in \mathrm{A}, x \in \mathrm{M};\tag{8}
\]

则用 \(\alpha x\) 记两边的公共值。

也可以说，在 M 上给定一个关于 E 和 F 的双模结构，等价于在 M 上给定一个 A-模结构，并且再给定两个 A-代数同态 \(s \mapsto h_s'\) （从 E 到 \(\mathrm{End}_{\mathbf{A}}(\mathbf{M})\) ）和 \(t \mapsto h_t''\) （从 F 到 \(\mathrm{End}_{\mathbf{A}}(\mathbf{M})\) ），使得 \(h_s'h_t'' = h_t''h_s'\) 对所有 \(s \in \mathbf{E}\) 和 \(t \in \mathbf{F}\) 成立。因此（第 2 号, 命题 5）典范地导出一个 A-代数同态 \(u \mapsto h_u\) （从 E \(\otimes_{\mathbf{A}} \mathbf{F}\) 到 \(\mathrm{End}_{\mathbf{A}}(\mathbf{M})\) ），使得 \(h_{s \otimes t} = h_s'h_t'' = h_t''h_s'\) 对 \(s \in \mathbf{E}\) 和 \(t \in \mathbf{F}\) 成立。换言之，这样就在 M 上定义了一个 \((\mathbf{E} \otimes_{\mathbf{A}} \mathbf{F})\) -模结构，称它为与给定的 E 和 F 上的双模结构相关联的结构，在此结构下

\[
(s \otimes t). x = s (t x) = t (s x) \quad \text { for } s \in \mathrm{E}, t \in \mathrm{F} \text { and } x \in \mathrm{M}.
\]

M 上给定的 E-模结构与 F-模结构可以由这个 \((\mathbf{E} \otimes_{\mathbf{A}} \mathbf{F})\) -模结构通过标量环的限制得到，它们分别对应两个典范同态 \(E \to E \otimes_{A} F\) 和 \(F \to E \otimes_{A} F\) .

反之，若在 M 上给定了一个 \((\mathbf{E} \otimes_{\mathbf{A}} \mathbf{F})\) -模结构，则借助典范同态 \(E \to E \otimes_{A} F\) 和 \(F \to E \otimes_{A} F\) 就得到 M 上的一个 E-模结构和一个 F-模结构，并且显然，带有这两个结构的 M 是代数 E 和 F 上的双模，且给定的 \((\mathbf{E} \otimes_{\mathbf{A}} \mathbf{F})\) -模结构与这个双模结构相关联。

这样就在 \((\mathbf{E} \otimes_{\mathbf{A}} \mathbf{F})\) -模与代数 E 和 F 上的双模之间建立了一个一一对应。显然，M 的每个子双模都是相伴 \((\mathbf{E} \otimes_{\mathbf{A}} \mathbf{F})\) -模结构的子模，反之亦然。对商、积、直和以及逆向极限与正向极限也有类似的结果。最后，若 \(\mathbf{M}'\) 是代数 \(\mathbf{E}\) 和 \(\mathbf{F}\) 上的另一个双模，且 \(f\colon \mathbf{M}\to \mathbf{M}'\) 是双模同态，则 \(f\) 也是 \((\mathbf{E}\otimes_{\mathbf{A}}\mathbf{F})\) -模同态，反之亦然。

对右双模结构，或者例如当存在左 E-模结构和右 F-模结构时，显然有相应的陈述；在这种情形我们说 (E, F)-双模，且给定这样一个结构等价于给定一个关于 E 和 F° 的左双模结构。

例。(1) 设 B 是一个 A-代数；环 B 典范地具有 (B, B)-双模结构（II, §1, 第 14 号, 例 1），且若 e 是 B 的单位元，则 \((\alpha e)x = x(\alpha e) = \alpha x\) 对所有 \(x \in B\) 及所有 \(\alpha \in A\) 成立；因此 B 可以看作代数 B 和 \(B^{\circ}\) （B 的反代数）上的左双模；从而与 B 上的 (B, B)-双模结构相关联，有一个 \((\mathbf{B} \otimes_{\mathbf{A}} \mathbf{B}^{\circ})\) -模结构，使得对 B 中的 b, x 和 \(b'\) ，

\[
(b \otimes b ^ {\prime}). x = b x b ^ {\prime}\tag{9}
\]

其中右侧是环B中的乘积。

\begin{enumerate}
\def\labelenumi{(\arabic{enumi})}
\setcounter{enumi}{1}
\tightlist
\item
  设 E 和 F 是两个 A-代数，e, \(e'\) 是它们各自的单位元，M 是一个 E-模，N 是一个 F-模；这些模结构在 M 上定义了关于环 A 和 E 的双模结构，在 N 上定义了关于环 A 和 F 的双模结构；由此导出张量积 \(M \otimes_{A} N\) 上关于环 E 和 F 的双模结构，定义为
\end{enumerate}

\[
x. (m \otimes n) = (x. m) \otimes n, \quad y. (m \otimes n) = m \otimes (y. n)
\]

对于 \(x \in \mathbf{E}, y \in \mathbf{F}, m \in \mathbf{M}, n \in \mathbf{N}\) （II, §3, 第 4 号）；还可看出条件 (8) 成立，因此上述双模结构伴随有一个 \((\mathbf{E} \otimes_{\mathbf{A}} \mathbf{F})\) -模结构（在 \(\mathbf{M} \otimes_{\mathbf{A}} \mathbf{N}\) 上），使得

\[
(x \otimes y). (m \otimes n) = (x. m) \otimes (y. n)\tag{10}
\]

对于 \(x \in \mathbf{E}, y \in \mathbf{F}, m \in \mathbf{M}, n \in \mathbf{N}\) .

特别地，取 \(\mathbf{M} = \mathbf{E}_s\) ， \(\mathbf{E}_s \otimes_{\mathbf{A}} \mathbf{N}\) 典范地具有一个 \((\mathbf{E} \otimes_{\mathbf{A}} \mathbf{F})\) -模结构；另一方面， \(\mathbf{E} \otimes_{\mathbf{A}} \mathbf{N}\) 典范地等同于 \(\mathbf{E} \otimes_{\mathbf{A}} (\mathbf{F}_d \otimes_{\mathbf{F}} \mathbf{N}) = (\mathbf{E} \otimes_{\mathbf{A}} \mathbf{F}) \otimes_{\mathbf{F}} \mathbf{N}\) ，其中 \(\mathbf{E} \otimes_{\mathbf{A}} \mathbf{F}\) 被视为带有由典范同态 \(v: \mathbf{F} \to \mathbf{E} \otimes_{\mathbf{A}} \mathbf{F}\) 定义的右 F-模结构；对 \(x, x'\) 在 \(\mathbf{E}, y \in \mathbf{F}, n \in \mathbf{N}\) ， \(x' \otimes n\) 于是等同于 \((x' \otimes e') \otimes n\) ，而 \((x \otimes y).(x' \otimes n') = (xx') \otimes (y.n)\) 等同于 \(((xx') \otimes y) \otimes n\) 。于是 \((\mathbf{E} \otimes_{\mathbf{A}} \mathbf{F})\) -模 \(\mathbf{E}_s \otimes_{\mathbf{A}} \mathbf{N}\) 等同于 \((\mathbf{E} \otimes_{\mathbf{A}} \mathbf{F})\) -模，即由 N 借助同态把标量扩张到 \(\mathbf{E} \otimes_{\mathbf{A}} \mathbf{F}\) 而得到的模，其中所用同态为 \(v\) （II, §5, 第 1 号）。 典范映射 \(n \mapsto e \otimes n\) （N 到 \(\mathbf{E}_s \otimes_{\mathbf{A}} \mathbf{N}\) 的）等同于典范映射 \(n \mapsto (e \otimes e') \otimes n\) （N 到 \((\mathbf{E} \otimes_{\mathbf{A}} \mathbf{F}) \otimes_{\mathbf{F}} \mathbf{N}\) 的）；已知这是一个 F-同态。

采用相同记号，设 P 是一个右 \((\mathbf{E} \otimes_{\mathbf{A}} \mathbf{F})\) -模；则存在一个典范的 Z-模同构

\[
\mathbf {P} \otimes_ {\mathbf {E} \otimes_ {\mathbf {A}} F} (\mathbf {E} _ {s} \otimes_ {\mathbf {A}} \mathbf {N}) \rightarrow \mathbf {P} \otimes_ {\mathbf {F}} \mathbf {N}\tag{11}
\]

其中右侧的 P 通过典范同态 v 被视为右 F-模。因为 P 典范地等同于 \(\mathrm{P}\otimes_{\mathrm{E}\otimes_{\mathrm{A}}\mathrm{F}}(\mathrm{E}\otimes_{\mathrm{A}}\mathrm{F})\) ，且 \((\mathrm{E}\otimes_{\mathrm{A}}\mathrm{F})\otimes_{\mathrm{F}}\mathrm{N}\) 等同于 \(\mathrm{E}\otimes_{\mathrm{A}}(\mathrm{F}\otimes_{\mathrm{F}}\mathrm{N})\) ，从而等同于 \(E\otimes_{A}N\) ，这就建立了所述的同构（II, §3, 第 8 号, 命题 8 及 II, §3, 第 4 号, 命题 4）。

以上所有内容都推广到多重模（II, §1, 第 14 号）。

